\chapter{Comparaison de toute la chaîne}\label{ch:comparaison}
\iterbanner{V1}

\section{Le protocole}
Toutes les méthodes passent par le même compteur et sont arrêtées dès que le front exact est entièrement simulé (budget maximal : 150 simulations pour la boucle). On répète chaque expérience sur 20 graines (10 pour QAOA, plus coûteux à simuler) et on rapporte la médiane et les quartiles de deux grandeurs :
\begin{itemize}
\item le nombre de simulations pour avoir simulé \textbf{tout le front exact} ;
\item le nombre de simulations pour atteindre \textbf{99~\% de l'hypervolume exact}, c'est-à-dire un front pratiquement aussi bon que le vrai.
\end{itemize}
Les méthodes comparées : tirage aléatoire sans remise (la référence naïve), NSGA-II ($N=16$, et $N=8$ pour le réglage optimal du chapitre~\ref{ch:nsga-res}), et la boucle QUBO avec trois solveurs : exact, recuit simulé, QAOA $p=2$. Le script est \code{optim/experiences\_pipeline.py} ; il tourne environ une heure, dont l'essentiel pour QAOA.

\section{Les courbes de convergence}
\begin{figure}[htbp]
\centering
% figure générée par tools/make_figures.py : convergence_front
\begin{tikzpicture}
\begin{groupplot}[group style={group size=2 by 1,horizontal sep=1.4cm},width=0.5\linewidth,height=6cm,
  xlabel={simulations distinctes},ylabel={fraction du front exact trouvée},xmin=0,xmax=300,ymin=0,ymax=1.02,
  title style={font=\small},tick label style={font=\scriptsize},
  legend style={draw=none,font=\scriptsize,at={(1.1,-0.25)},anchor=north},legend columns=3,legend cell align=left]
\nextgroupplot[title={S1 : approche de face (7 profils sur le front)}]
\addlegendimage{cmuted,very thick}\addlegendentry{tirage aléatoire}
\addlegendimage{corange,very thick}\addlegendentry{NSGA-II}
\addlegendimage{cblue,very thick}\addlegendentry{boucle QUBO, solveur exact}
\addlegendimage{cgreen,very thick}\addlegendentry{boucle QUBO, recuit simulé}
\addlegendimage{cviolet,very thick}\addlegendentry{boucle QUBO, QAOA $p=2$}
\addplot[cmuted,very thick,forget plot] coordinates {(0,0) (2,0) (4,0) (6,0) (8,0) (10,0) (12,0) (14,0) (16,0) (18,0) (20,0) (22,0) (24,0) (26,0) (28,0) (30,0) (32,0) (34,0) (36,0) (38,0) (40,0) (42,0) (44,0) (46,0) (48,0) (50,0) (52,0) (54,0) (56,0.071429) (58,0.071429) (60,0.071429) (62,0.071429) (64,0.071429) (66,0.071429) (68,0.14286) (70,0.14286) (72,0.14286) (74,0.14286) (76,0.14286) (78,0.14286) (80,0.14286) (82,0.14286) (84,0.14286) (86,0.14286) (88,0.14286) (90,0.14286) (92,0.14286) (94,0.14286) (96,0.14286) (98,0.14286) (100,0.14286) (102,0.14286) (104,0.14286) (106,0.14286) (108,0.21429) (110,0.21429) (112,0.21429) (114,0.21429) (116,0.21429) (118,0.21429) (120,0.21429) (122,0.21429) (124,0.28571) (126,0.28571) (128,0.28571) (130,0.28571) (132,0.28571) (134,0.28571) (136,0.28571) (138,0.28571) (140,0.28571) (142,0.28571) (144,0.28571) (146,0.28571) (148,0.28571) (150,0.28571) (152,0.28571) (154,0.28571) (156,0.28571) (158,0.28571) (160,0.28571) (162,0.28571) (164,0.28571) (166,0.28571) (168,0.28571) (170,0.28571) (172,0.28571) (174,0.28571) (176,0.28571) (178,0.28571) (180,0.28571) (182,0.28571) (184,0.28571) (186,0.28571) (188,0.28571) (190,0.28571) (192,0.35714) (194,0.35714) (196,0.35714) (198,0.35714) (200,0.35714) (202,0.35714) (204,0.35714) (206,0.35714) (208,0.42857) (210,0.42857) (212,0.42857) (214,0.42857) (216,0.42857) (218,0.42857) (220,0.42857) (222,0.42857) (224,0.42857) (226,0.42857) (228,0.42857) (230,0.42857) (232,0.42857) (234,0.42857) (236,0.42857) (238,0.42857) (240,0.42857) (242,0.42857) (244,0.42857) (246,0.42857) (248,0.42857) (250,0.42857) (252,0.42857) (254,0.42857) (256,0.42857) (258,0.42857) (260,0.57143) (262,0.57143) (264,0.57143) (266,0.57143) (268,0.57143) (270,0.57143) (272,0.57143) (274,0.57143) (276,0.57143) (278,0.57143) (280,0.57143) (282,0.57143) (284,0.57143) (286,0.57143) (288,0.57143) (290,0.57143) (292,0.57143) (294,0.57143) (296,0.57143) (298,0.57143) (300,0.57143)};
\addplot[cmuted,thin,dashed,forget plot] coordinates {(0,0) (2,0) (4,0) (6,0) (8,0) (10,0) (12,0) (14,0) (16,0) (18,0) (20,0) (22,0) (24,0) (26,0) (28,0) (30,0) (32,0) (34,0) (36,0) (38,0) (40,0) (42,0) (44,0) (46,0) (48,0) (50,0) (52,0) (54,0) (56,0) (58,0) (60,0) (62,0) (64,0) (66,0) (68,0) (70,0) (72,0) (74,0) (76,0) (78,0) (80,0) (82,0) (84,0) (86,0) (88,0) (90,0) (92,0) (94,0) (96,0.10714) (98,0.10714) (100,0.14286) (102,0.14286) (104,0.14286) (106,0.14286) (108,0.14286) (110,0.14286) (112,0.14286) (114,0.14286) (116,0.14286) (118,0.14286) (120,0.14286) (122,0.14286) (124,0.14286) (126,0.14286) (128,0.14286) (130,0.14286) (132,0.14286) (134,0.14286) (136,0.14286) (138,0.14286) (140,0.14286) (142,0.14286) (144,0.14286) (146,0.14286) (148,0.14286) (150,0.14286) (152,0.14286) (154,0.14286) (156,0.14286) (158,0.14286) (160,0.14286) (162,0.14286) (164,0.14286) (166,0.14286) (168,0.14286) (170,0.14286) (172,0.14286) (174,0.14286) (176,0.14286) (178,0.14286) (180,0.14286) (182,0.14286) (184,0.25) (186,0.25) (188,0.25) (190,0.25) (192,0.25) (194,0.25) (196,0.28571) (198,0.28571) (200,0.28571) (202,0.28571) (204,0.28571) (206,0.28571) (208,0.28571) (210,0.28571) (212,0.28571) (214,0.28571) (216,0.28571) (218,0.28571) (220,0.28571) (222,0.28571) (224,0.28571) (226,0.28571) (228,0.28571) (230,0.28571) (232,0.28571) (234,0.28571) (236,0.28571) (238,0.28571) (240,0.28571) (242,0.28571) (244,0.28571) (246,0.28571) (248,0.28571) (250,0.28571) (252,0.28571) (254,0.28571) (256,0.28571) (258,0.28571) (260,0.28571) (262,0.28571) (264,0.28571) (266,0.28571) (268,0.28571) (270,0.28571) (272,0.39286) (274,0.39286) (276,0.39286) (278,0.39286) (280,0.42857) (282,0.42857) (284,0.42857) (286,0.42857) (288,0.42857) (290,0.42857) (292,0.42857) (294,0.42857) (296,0.42857) (298,0.42857) (300,0.42857)};
\addplot[corange,very thick,forget plot] coordinates {(0,0) (2,0) (4,0) (6,0) (8,0) (10,0) (12,0) (14,0) (16,0) (18,0) (20,0) (22,0) (24,0) (26,0) (28,0) (30,0) (32,0) (34,0) (36,0) (38,0) (40,0) (42,0) (44,0) (46,0) (48,0.14286) (50,0.14286) (52,0.14286) (54,0.14286) (56,0.14286) (58,0.14286) (60,0.14286) (62,0.28571) (64,0.35714) (66,0.35714) (68,0.35714) (70,0.35714) (72,0.35714) (74,0.42857) (76,0.42857) (78,0.5) (80,0.57143) (82,0.57143) (84,0.57143) (86,0.57143) (88,0.71429) (90,0.78571) (92,0.85714) (94,0.85714) (96,0.85714) (98,0.85714) (100,0.85714) (102,0.85714) (104,0.92857) (106,0.92857) (108,0.92857) (110,1) (112,1) (114,1) (116,1) (118,1) (120,1) (122,1) (124,1) (126,1) (128,1) (130,1) (132,1) (134,1) (136,1) (138,1) (140,1) (142,1) (144,1) (146,1) (148,1) (150,1) (152,1) (154,1) (156,1) (158,1) (160,1) (162,1) (164,1) (166,1) (168,1) (170,1) (172,1) (174,1) (176,1) (178,1) (180,1) (182,1) (184,1) (186,1) (188,1) (190,1) (192,1) (194,1) (196,1) (198,1) (200,1) (202,1) (204,1) (206,1) (208,1) (210,1) (212,1) (214,1) (216,1) (218,1) (220,1) (222,1) (224,1) (226,1) (228,1) (230,1) (232,1) (234,1) (236,1) (238,1) (240,1) (242,1) (244,1) (246,1) (248,1) (250,1) (252,1) (254,1) (256,1) (258,1) (260,1) (262,1) (264,1) (266,1) (268,1) (270,1) (272,1) (274,1) (276,1) (278,1) (280,1) (282,1) (284,1) (286,1) (288,1) (290,1) (292,1) (294,1) (296,1) (298,1) (300,1)};
\addplot[corange,thin,dashed,forget plot] coordinates {(0,0) (2,0) (4,0) (6,0) (8,0) (10,0) (12,0) (14,0) (16,0) (18,0) (20,0) (22,0) (24,0) (26,0) (28,0) (30,0) (32,0) (34,0) (36,0) (38,0) (40,0) (42,0) (44,0) (46,0) (48,0.10714) (50,0.10714) (52,0.10714) (54,0.10714) (56,0.10714) (58,0.10714) (60,0.10714) (62,0.10714) (64,0.28571) (66,0.28571) (68,0.28571) (70,0.28571) (72,0.28571) (74,0.28571) (76,0.39286) (78,0.42857) (80,0.42857) (82,0.42857) (84,0.42857) (86,0.42857) (88,0.57143) (90,0.57143) (92,0.57143) (94,0.67857) (96,0.85714) (98,0.85714) (100,0.85714) (102,0.85714) (104,0.85714) (106,0.85714) (108,0.85714) (110,0.85714) (112,0.85714) (114,1) (116,1) (118,1) (120,1) (122,1) (124,1) (126,1) (128,1) (130,1) (132,1) (134,1) (136,1) (138,1) (140,1) (142,1) (144,1) (146,1) (148,1) (150,1) (152,1) (154,1) (156,1) (158,1) (160,1) (162,1) (164,1) (166,1) (168,1) (170,1) (172,1) (174,1) (176,1) (178,1) (180,1) (182,1) (184,1) (186,1) (188,1) (190,1) (192,1) (194,1) (196,1) (198,1) (200,1) (202,1) (204,1) (206,1) (208,1) (210,1) (212,1) (214,1) (216,1) (218,1) (220,1) (222,1) (224,1) (226,1) (228,1) (230,1) (232,1) (234,1) (236,1) (238,1) (240,1) (242,1) (244,1) (246,1) (248,1) (250,1) (252,1) (254,1) (256,1) (258,1) (260,1) (262,1) (264,1) (266,1) (268,1) (270,1) (272,1) (274,1) (276,1) (278,1) (280,1) (282,1) (284,1) (286,1) (288,1) (290,1) (292,1) (294,1) (296,1) (298,1) (300,1)};
\addplot[cblue,very thick,forget plot] coordinates {(0,0) (2,0) (4,0) (6,0) (8,0) (10,0) (12,0) (14,0) (16,0.28571) (18,0.28571) (20,0.42857) (22,0.42857) (24,0.42857) (26,0.78571) (28,0.78571) (30,0.85714) (32,0.85714) (34,0.85714) (36,1) (38,1) (40,1) (42,1) (44,1) (46,1) (48,1) (50,1) (52,1) (54,1) (56,1) (58,1) (60,1) (62,1) (64,1) (66,1) (68,1) (70,1) (72,1) (74,1) (76,1) (78,1) (80,1) (82,1) (84,1) (86,1) (88,1) (90,1) (92,1) (94,1) (96,1) (98,1) (100,1) (102,1) (104,1) (106,1) (108,1) (110,1) (112,1) (114,1) (116,1) (118,1) (120,1) (122,1) (124,1) (126,1) (128,1) (130,1) (132,1) (134,1) (136,1) (138,1) (140,1) (142,1) (144,1) (146,1) (148,1) (150,1) (152,1) (154,1) (156,1) (158,1) (160,1) (162,1) (164,1) (166,1) (168,1) (170,1) (172,1) (174,1) (176,1) (178,1) (180,1) (182,1) (184,1) (186,1) (188,1) (190,1) (192,1) (194,1) (196,1) (198,1) (200,1) (202,1) (204,1) (206,1) (208,1) (210,1) (212,1) (214,1) (216,1) (218,1) (220,1) (222,1) (224,1) (226,1) (228,1) (230,1) (232,1) (234,1) (236,1) (238,1) (240,1) (242,1) (244,1) (246,1) (248,1) (250,1) (252,1) (254,1) (256,1) (258,1) (260,1) (262,1) (264,1) (266,1) (268,1) (270,1) (272,1) (274,1) (276,1) (278,1) (280,1) (282,1) (284,1) (286,1) (288,1) (290,1) (292,1) (294,1) (296,1) (298,1) (300,1)};
\addplot[cblue,thin,dashed,forget plot] coordinates {(0,0) (2,0) (4,0) (6,0) (8,0) (10,0) (12,0) (14,0) (16,0.10714) (18,0.10714) (20,0.42857) (22,0.42857) (24,0.42857) (26,0.57143) (28,0.57143) (30,0.85714) (32,0.85714) (34,0.85714) (36,0.85714) (38,0.85714) (40,1) (42,1) (44,1) (46,1) (48,1) (50,1) (52,1) (54,1) (56,1) (58,1) (60,1) (62,1) (64,1) (66,1) (68,1) (70,1) (72,1) (74,1) (76,1) (78,1) (80,1) (82,1) (84,1) (86,1) (88,1) (90,1) (92,1) (94,1) (96,1) (98,1) (100,1) (102,1) (104,1) (106,1) (108,1) (110,1) (112,1) (114,1) (116,1) (118,1) (120,1) (122,1) (124,1) (126,1) (128,1) (130,1) (132,1) (134,1) (136,1) (138,1) (140,1) (142,1) (144,1) (146,1) (148,1) (150,1) (152,1) (154,1) (156,1) (158,1) (160,1) (162,1) (164,1) (166,1) (168,1) (170,1) (172,1) (174,1) (176,1) (178,1) (180,1) (182,1) (184,1) (186,1) (188,1) (190,1) (192,1) (194,1) (196,1) (198,1) (200,1) (202,1) (204,1) (206,1) (208,1) (210,1) (212,1) (214,1) (216,1) (218,1) (220,1) (222,1) (224,1) (226,1) (228,1) (230,1) (232,1) (234,1) (236,1) (238,1) (240,1) (242,1) (244,1) (246,1) (248,1) (250,1) (252,1) (254,1) (256,1) (258,1) (260,1) (262,1) (264,1) (266,1) (268,1) (270,1) (272,1) (274,1) (276,1) (278,1) (280,1) (282,1) (284,1) (286,1) (288,1) (290,1) (292,1) (294,1) (296,1) (298,1) (300,1)};
\addplot[cgreen,very thick,forget plot] coordinates {(0,0) (2,0) (4,0) (6,0) (8,0) (10,0) (12,0) (14,0) (16,0.28571) (18,0.28571) (20,0.42857) (22,0.42857) (24,0.42857) (26,0.71429) (28,0.71429) (30,0.85714) (32,0.85714) (34,0.85714) (36,0.85714) (38,0.85714) (40,1) (42,1) (44,1) (46,1) (48,1) (50,1) (52,1) (54,1) (56,1) (58,1) (60,1) (62,1) (64,1) (66,1) (68,1) (70,1) (72,1) (74,1) (76,1) (78,1) (80,1) (82,1) (84,1) (86,1) (88,1) (90,1) (92,1) (94,1) (96,1) (98,1) (100,1) (102,1) (104,1) (106,1) (108,1) (110,1) (112,1) (114,1) (116,1) (118,1) (120,1) (122,1) (124,1) (126,1) (128,1) (130,1) (132,1) (134,1) (136,1) (138,1) (140,1) (142,1) (144,1) (146,1) (148,1) (150,1) (152,1) (154,1) (156,1) (158,1) (160,1) (162,1) (164,1) (166,1) (168,1) (170,1) (172,1) (174,1) (176,1) (178,1) (180,1) (182,1) (184,1) (186,1) (188,1) (190,1) (192,1) (194,1) (196,1) (198,1) (200,1) (202,1) (204,1) (206,1) (208,1) (210,1) (212,1) (214,1) (216,1) (218,1) (220,1) (222,1) (224,1) (226,1) (228,1) (230,1) (232,1) (234,1) (236,1) (238,1) (240,1) (242,1) (244,1) (246,1) (248,1) (250,1) (252,1) (254,1) (256,1) (258,1) (260,1) (262,1) (264,1) (266,1) (268,1) (270,1) (272,1) (274,1) (276,1) (278,1) (280,1) (282,1) (284,1) (286,1) (288,1) (290,1) (292,1) (294,1) (296,1) (298,1) (300,1)};
\addplot[cgreen,thin,dashed,forget plot] coordinates {(0,0) (2,0) (4,0) (6,0) (8,0) (10,0) (12,0) (14,0) (16,0.14286) (18,0.14286) (20,0.28571) (22,0.28571) (24,0.28571) (26,0.53571) (28,0.53571) (30,0.67857) (32,0.67857) (34,0.67857) (36,0.71429) (38,0.71429) (40,0.85714) (42,0.85714) (44,0.85714) (46,0.96429) (48,0.96429) (50,1) (52,1) (54,1) (56,1) (58,1) (60,1) (62,1) (64,1) (66,1) (68,1) (70,1) (72,1) (74,1) (76,1) (78,1) (80,1) (82,1) (84,1) (86,1) (88,1) (90,1) (92,1) (94,1) (96,1) (98,1) (100,1) (102,1) (104,1) (106,1) (108,1) (110,1) (112,1) (114,1) (116,1) (118,1) (120,1) (122,1) (124,1) (126,1) (128,1) (130,1) (132,1) (134,1) (136,1) (138,1) (140,1) (142,1) (144,1) (146,1) (148,1) (150,1) (152,1) (154,1) (156,1) (158,1) (160,1) (162,1) (164,1) (166,1) (168,1) (170,1) (172,1) (174,1) (176,1) (178,1) (180,1) (182,1) (184,1) (186,1) (188,1) (190,1) (192,1) (194,1) (196,1) (198,1) (200,1) (202,1) (204,1) (206,1) (208,1) (210,1) (212,1) (214,1) (216,1) (218,1) (220,1) (222,1) (224,1) (226,1) (228,1) (230,1) (232,1) (234,1) (236,1) (238,1) (240,1) (242,1) (244,1) (246,1) (248,1) (250,1) (252,1) (254,1) (256,1) (258,1) (260,1) (262,1) (264,1) (266,1) (268,1) (270,1) (272,1) (274,1) (276,1) (278,1) (280,1) (282,1) (284,1) (286,1) (288,1) (290,1) (292,1) (294,1) (296,1) (298,1) (300,1)};
\addplot[cviolet,very thick,forget plot] coordinates {(0,0) (2,0) (4,0) (6,0) (8,0) (10,0) (12,0) (14,0) (16,0.21429) (18,0.21429) (20,0.42857) (22,0.42857) (24,0.42857) (26,0.71429) (28,0.71429) (30,0.85714) (32,0.85714) (34,0.85714) (36,1) (38,1) (40,1) (42,1) (44,1) (46,1) (48,1) (50,1) (52,1) (54,1) (56,1) (58,1) (60,1) (62,1) (64,1) (66,1) (68,1) (70,1) (72,1) (74,1) (76,1) (78,1) (80,1) (82,1) (84,1) (86,1) (88,1) (90,1) (92,1) (94,1) (96,1) (98,1) (100,1) (102,1) (104,1) (106,1) (108,1) (110,1) (112,1) (114,1) (116,1) (118,1) (120,1) (122,1) (124,1) (126,1) (128,1) (130,1) (132,1) (134,1) (136,1) (138,1) (140,1) (142,1) (144,1) (146,1) (148,1) (150,1) (152,1) (154,1) (156,1) (158,1) (160,1) (162,1) (164,1) (166,1) (168,1) (170,1) (172,1) (174,1) (176,1) (178,1) (180,1) (182,1) (184,1) (186,1) (188,1) (190,1) (192,1) (194,1) (196,1) (198,1) (200,1) (202,1) (204,1) (206,1) (208,1) (210,1) (212,1) (214,1) (216,1) (218,1) (220,1) (222,1) (224,1) (226,1) (228,1) (230,1) (232,1) (234,1) (236,1) (238,1) (240,1) (242,1) (244,1) (246,1) (248,1) (250,1) (252,1) (254,1) (256,1) (258,1) (260,1) (262,1) (264,1) (266,1) (268,1) (270,1) (272,1) (274,1) (276,1) (278,1) (280,1) (282,1) (284,1) (286,1) (288,1) (290,1) (292,1) (294,1) (296,1) (298,1) (300,1)};
\addplot[cviolet,thin,dashed,forget plot] coordinates {(0,0) (2,0) (4,0) (6,0) (8,0) (10,0) (12,0) (14,0) (16,0.14286) (18,0.14286) (20,0.42857) (22,0.42857) (24,0.42857) (26,0.60714) (28,0.60714) (30,0.85714) (32,0.85714) (34,0.85714) (36,1) (38,1) (40,1) (42,1) (44,1) (46,1) (48,1) (50,1) (52,1) (54,1) (56,1) (58,1) (60,1) (62,1) (64,1) (66,1) (68,1) (70,1) (72,1) (74,1) (76,1) (78,1) (80,1) (82,1) (84,1) (86,1) (88,1) (90,1) (92,1) (94,1) (96,1) (98,1) (100,1) (102,1) (104,1) (106,1) (108,1) (110,1) (112,1) (114,1) (116,1) (118,1) (120,1) (122,1) (124,1) (126,1) (128,1) (130,1) (132,1) (134,1) (136,1) (138,1) (140,1) (142,1) (144,1) (146,1) (148,1) (150,1) (152,1) (154,1) (156,1) (158,1) (160,1) (162,1) (164,1) (166,1) (168,1) (170,1) (172,1) (174,1) (176,1) (178,1) (180,1) (182,1) (184,1) (186,1) (188,1) (190,1) (192,1) (194,1) (196,1) (198,1) (200,1) (202,1) (204,1) (206,1) (208,1) (210,1) (212,1) (214,1) (216,1) (218,1) (220,1) (222,1) (224,1) (226,1) (228,1) (230,1) (232,1) (234,1) (236,1) (238,1) (240,1) (242,1) (244,1) (246,1) (248,1) (250,1) (252,1) (254,1) (256,1) (258,1) (260,1) (262,1) (264,1) (266,1) (268,1) (270,1) (272,1) (274,1) (276,1) (278,1) (280,1) (282,1) (284,1) (286,1) (288,1) (290,1) (292,1) (294,1) (296,1) (298,1) (300,1)};
\nextgroupplot[title={S2 : radar en contrebas (15 profils sur le front)}]
\addplot[cmuted,very thick,forget plot] coordinates {(0,0) (2,0) (4,0) (6,0) (8,0) (10,0) (12,0) (14,0) (16,0) (18,0) (20,0) (22,0) (24,0) (26,0) (28,0) (30,0) (32,0.033333) (34,0.033333) (36,0.033333) (38,0.033333) (40,0.033333) (42,0.033333) (44,0.066667) (46,0.066667) (48,0.066667) (50,0.066667) (52,0.066667) (54,0.066667) (56,0.066667) (58,0.066667) (60,0.066667) (62,0.066667) (64,0.066667) (66,0.066667) (68,0.066667) (70,0.066667) (72,0.1) (74,0.1) (76,0.13333) (78,0.13333) (80,0.13333) (82,0.13333) (84,0.13333) (86,0.13333) (88,0.16667) (90,0.16667) (92,0.2) (94,0.2) (96,0.2) (98,0.2) (100,0.2) (102,0.2) (104,0.2) (106,0.2) (108,0.2) (110,0.2) (112,0.2) (114,0.2) (116,0.26667) (118,0.26667) (120,0.26667) (122,0.26667) (124,0.26667) (126,0.26667) (128,0.26667) (130,0.26667) (132,0.26667) (134,0.26667) (136,0.26667) (138,0.26667) (140,0.26667) (142,0.26667) (144,0.26667) (146,0.26667) (148,0.3) (150,0.3) (152,0.3) (154,0.3) (156,0.3) (158,0.3) (160,0.3) (162,0.3) (164,0.3) (166,0.3) (168,0.3) (170,0.3) (172,0.33333) (174,0.33333) (176,0.33333) (178,0.33333) (180,0.33333) (182,0.33333) (184,0.33333) (186,0.33333) (188,0.33333) (190,0.33333) (192,0.33333) (194,0.33333) (196,0.33333) (198,0.33333) (200,0.36667) (202,0.36667) (204,0.36667) (206,0.36667) (208,0.4) (210,0.4) (212,0.4) (214,0.4) (216,0.4) (218,0.4) (220,0.4) (222,0.4) (224,0.43333) (226,0.43333) (228,0.43333) (230,0.43333) (232,0.46667) (234,0.46667) (236,0.46667) (238,0.46667) (240,0.46667) (242,0.46667) (244,0.46667) (246,0.46667) (248,0.46667) (250,0.46667) (252,0.5) (254,0.5) (256,0.53333) (258,0.53333) (260,0.53333) (262,0.53333) (264,0.53333) (266,0.53333) (268,0.53333) (270,0.53333) (272,0.56667) (274,0.56667) (276,0.56667) (278,0.56667) (280,0.56667) (282,0.56667) (284,0.56667) (286,0.56667) (288,0.56667) (290,0.56667) (292,0.6) (294,0.6) (296,0.6) (298,0.6) (300,0.6)};
\addplot[cmuted,thin,dashed,forget plot] coordinates {(0,0) (2,0) (4,0) (6,0) (8,0) (10,0) (12,0) (14,0) (16,0) (18,0) (20,0) (22,0) (24,0) (26,0) (28,0) (30,0) (32,0) (34,0) (36,0) (38,0) (40,0) (42,0) (44,0) (46,0) (48,0) (50,0) (52,0) (54,0) (56,0) (58,0) (60,0) (62,0) (64,0) (66,0) (68,0.05) (70,0.05) (72,0.066667) (74,0.066667) (76,0.066667) (78,0.066667) (80,0.066667) (82,0.066667) (84,0.066667) (86,0.066667) (88,0.066667) (90,0.066667) (92,0.066667) (94,0.066667) (96,0.11667) (98,0.11667) (100,0.13333) (102,0.13333) (104,0.13333) (106,0.13333) (108,0.13333) (110,0.13333) (112,0.13333) (114,0.13333) (116,0.13333) (118,0.13333) (120,0.13333) (122,0.13333) (124,0.13333) (126,0.13333) (128,0.13333) (130,0.13333) (132,0.13333) (134,0.13333) (136,0.13333) (138,0.13333) (140,0.13333) (142,0.13333) (144,0.18333) (146,0.18333) (148,0.18333) (150,0.18333) (152,0.18333) (154,0.18333) (156,0.18333) (158,0.18333) (160,0.18333) (162,0.18333) (164,0.2) (166,0.2) (168,0.25) (170,0.25) (172,0.25) (174,0.25) (176,0.25) (178,0.25) (180,0.26667) (182,0.26667) (184,0.26667) (186,0.26667) (188,0.26667) (190,0.26667) (192,0.26667) (194,0.26667) (196,0.26667) (198,0.26667) (200,0.26667) (202,0.26667) (204,0.26667) (206,0.26667) (208,0.26667) (210,0.26667) (212,0.26667) (214,0.26667) (216,0.26667) (218,0.26667) (220,0.26667) (222,0.26667) (224,0.31667) (226,0.31667) (228,0.31667) (230,0.31667) (232,0.31667) (234,0.31667) (236,0.38333) (238,0.38333) (240,0.38333) (242,0.38333) (244,0.38333) (246,0.38333) (248,0.38333) (250,0.38333) (252,0.38333) (254,0.38333) (256,0.4) (258,0.4) (260,0.45) (262,0.45) (264,0.45) (266,0.45) (268,0.45) (270,0.45) (272,0.45) (274,0.45) (276,0.46667) (278,0.46667) (280,0.46667) (282,0.46667) (284,0.46667) (286,0.46667) (288,0.51667) (290,0.51667) (292,0.51667) (294,0.51667) (296,0.51667) (298,0.51667) (300,0.51667)};
\addplot[corange,very thick,forget plot] coordinates {(0,0) (2,0) (4,0) (6,0) (8,0) (10,0) (12,0) (14,0) (16,0.066667) (18,0.066667) (20,0.066667) (22,0.066667) (24,0.066667) (26,0.066667) (28,0.066667) (30,0.066667) (32,0.066667) (34,0.066667) (36,0.066667) (38,0.066667) (40,0.066667) (42,0.066667) (44,0.066667) (46,0.066667) (48,0.2) (50,0.2) (52,0.2) (54,0.2) (56,0.2) (58,0.2) (60,0.2) (62,0.23333) (64,0.33333) (66,0.33333) (68,0.33333) (70,0.33333) (72,0.33333) (74,0.36667) (76,0.43333) (78,0.53333) (80,0.6) (82,0.6) (84,0.6) (86,0.63333) (88,0.7) (90,0.7) (92,0.8) (94,0.83333) (96,0.83333) (98,0.93333) (100,0.93333) (102,0.93333) (104,0.93333) (106,0.96667) (108,1) (110,1) (112,1) (114,1) (116,1) (118,1) (120,1) (122,1) (124,1) (126,1) (128,1) (130,1) (132,1) (134,1) (136,1) (138,1) (140,1) (142,1) (144,1) (146,1) (148,1) (150,1) (152,1) (154,1) (156,1) (158,1) (160,1) (162,1) (164,1) (166,1) (168,1) (170,1) (172,1) (174,1) (176,1) (178,1) (180,1) (182,1) (184,1) (186,1) (188,1) (190,1) (192,1) (194,1) (196,1) (198,1) (200,1) (202,1) (204,1) (206,1) (208,1) (210,1) (212,1) (214,1) (216,1) (218,1) (220,1) (222,1) (224,1) (226,1) (228,1) (230,1) (232,1) (234,1) (236,1) (238,1) (240,1) (242,1) (244,1) (246,1) (248,1) (250,1) (252,1) (254,1) (256,1) (258,1) (260,1) (262,1) (264,1) (266,1) (268,1) (270,1) (272,1) (274,1) (276,1) (278,1) (280,1) (282,1) (284,1) (286,1) (288,1) (290,1) (292,1) (294,1) (296,1) (298,1) (300,1)};
\addplot[corange,thin,dashed,forget plot] coordinates {(0,0) (2,0) (4,0) (6,0) (8,0) (10,0) (12,0) (14,0) (16,0) (18,0) (20,0) (22,0) (24,0) (26,0) (28,0) (30,0) (32,0) (34,0) (36,0) (38,0) (40,0) (42,0) (44,0) (46,0) (48,0.11667) (50,0.11667) (52,0.11667) (54,0.11667) (56,0.11667) (58,0.11667) (60,0.13333) (62,0.13333) (64,0.2) (66,0.2) (68,0.2) (70,0.2) (72,0.2) (74,0.3) (76,0.3) (78,0.38333) (80,0.4) (82,0.4) (84,0.4) (86,0.45) (88,0.56667) (90,0.56667) (92,0.6) (94,0.6) (96,0.7) (98,0.78333) (100,0.78333) (102,0.85) (104,0.85) (106,0.9) (108,0.9) (110,0.93333) (112,0.98333) (114,1) (116,1) (118,1) (120,1) (122,1) (124,1) (126,1) (128,1) (130,1) (132,1) (134,1) (136,1) (138,1) (140,1) (142,1) (144,1) (146,1) (148,1) (150,1) (152,1) (154,1) (156,1) (158,1) (160,1) (162,1) (164,1) (166,1) (168,1) (170,1) (172,1) (174,1) (176,1) (178,1) (180,1) (182,1) (184,1) (186,1) (188,1) (190,1) (192,1) (194,1) (196,1) (198,1) (200,1) (202,1) (204,1) (206,1) (208,1) (210,1) (212,1) (214,1) (216,1) (218,1) (220,1) (222,1) (224,1) (226,1) (228,1) (230,1) (232,1) (234,1) (236,1) (238,1) (240,1) (242,1) (244,1) (246,1) (248,1) (250,1) (252,1) (254,1) (256,1) (258,1) (260,1) (262,1) (264,1) (266,1) (268,1) (270,1) (272,1) (274,1) (276,1) (278,1) (280,1) (282,1) (284,1) (286,1) (288,1) (290,1) (292,1) (294,1) (296,1) (298,1) (300,1)};
\addplot[cblue,very thick,forget plot] coordinates {(0,0) (2,0) (4,0) (6,0) (8,0) (10,0) (12,0) (14,0) (16,0.066667) (18,0.066667) (20,0.2) (22,0.2) (24,0.2) (26,0.33333) (28,0.33333) (30,0.4) (32,0.4) (34,0.4) (36,0.53333) (38,0.53333) (40,0.63333) (42,0.63333) (44,0.63333) (46,0.76667) (48,0.76667) (50,0.83333) (52,0.83333) (54,0.83333) (56,0.86667) (58,0.86667) (60,0.93333) (62,0.93333) (64,0.93333) (66,0.96667) (68,0.96667) (70,1) (72,1) (74,1) (76,1) (78,1) (80,1) (82,1) (84,1) (86,1) (88,1) (90,1) (92,1) (94,1) (96,1) (98,1) (100,1) (102,1) (104,1) (106,1) (108,1) (110,1) (112,1) (114,1) (116,1) (118,1) (120,1) (122,1) (124,1) (126,1) (128,1) (130,1) (132,1) (134,1) (136,1) (138,1) (140,1) (142,1) (144,1) (146,1) (148,1) (150,1) (152,1) (154,1) (156,1) (158,1) (160,1) (162,1) (164,1) (166,1) (168,1) (170,1) (172,1) (174,1) (176,1) (178,1) (180,1) (182,1) (184,1) (186,1) (188,1) (190,1) (192,1) (194,1) (196,1) (198,1) (200,1) (202,1) (204,1) (206,1) (208,1) (210,1) (212,1) (214,1) (216,1) (218,1) (220,1) (222,1) (224,1) (226,1) (228,1) (230,1) (232,1) (234,1) (236,1) (238,1) (240,1) (242,1) (244,1) (246,1) (248,1) (250,1) (252,1) (254,1) (256,1) (258,1) (260,1) (262,1) (264,1) (266,1) (268,1) (270,1) (272,1) (274,1) (276,1) (278,1) (280,1) (282,1) (284,1) (286,1) (288,1) (290,1) (292,1) (294,1) (296,1) (298,1) (300,1)};
\addplot[cblue,thin,dashed,forget plot] coordinates {(0,0) (2,0) (4,0) (6,0) (8,0) (10,0) (12,0) (14,0) (16,0.066667) (18,0.066667) (20,0.2) (22,0.2) (24,0.2) (26,0.26667) (28,0.26667) (30,0.4) (32,0.4) (34,0.4) (36,0.46667) (38,0.46667) (40,0.58333) (42,0.58333) (44,0.58333) (46,0.65) (48,0.65) (50,0.73333) (52,0.73333) (54,0.73333) (56,0.78333) (58,0.78333) (60,0.85) (62,0.85) (64,0.85) (66,0.86667) (68,0.86667) (70,0.86667) (72,0.86667) (74,0.86667) (76,0.91667) (78,0.91667) (80,0.93333) (82,0.93333) (84,0.93333) (86,1) (88,1) (90,1) (92,1) (94,1) (96,1) (98,1) (100,1) (102,1) (104,1) (106,1) (108,1) (110,1) (112,1) (114,1) (116,1) (118,1) (120,1) (122,1) (124,1) (126,1) (128,1) (130,1) (132,1) (134,1) (136,1) (138,1) (140,1) (142,1) (144,1) (146,1) (148,1) (150,1) (152,1) (154,1) (156,1) (158,1) (160,1) (162,1) (164,1) (166,1) (168,1) (170,1) (172,1) (174,1) (176,1) (178,1) (180,1) (182,1) (184,1) (186,1) (188,1) (190,1) (192,1) (194,1) (196,1) (198,1) (200,1) (202,1) (204,1) (206,1) (208,1) (210,1) (212,1) (214,1) (216,1) (218,1) (220,1) (222,1) (224,1) (226,1) (228,1) (230,1) (232,1) (234,1) (236,1) (238,1) (240,1) (242,1) (244,1) (246,1) (248,1) (250,1) (252,1) (254,1) (256,1) (258,1) (260,1) (262,1) (264,1) (266,1) (268,1) (270,1) (272,1) (274,1) (276,1) (278,1) (280,1) (282,1) (284,1) (286,1) (288,1) (290,1) (292,1) (294,1) (296,1) (298,1) (300,1)};
\addplot[cgreen,very thick,forget plot] coordinates {(0,0) (2,0) (4,0) (6,0) (8,0) (10,0) (12,0) (14,0) (16,0.066667) (18,0.066667) (20,0.2) (22,0.2) (24,0.2) (26,0.26667) (28,0.26667) (30,0.33333) (32,0.33333) (34,0.33333) (36,0.43333) (38,0.43333) (40,0.53333) (42,0.53333) (44,0.53333) (46,0.6) (48,0.6) (50,0.66667) (52,0.66667) (54,0.66667) (56,0.73333) (58,0.73333) (60,0.8) (62,0.8) (64,0.8) (66,0.83333) (68,0.83333) (70,0.86667) (72,0.86667) (74,0.86667) (76,0.9) (78,0.9) (80,0.93333) (82,0.93333) (84,0.93333) (86,1) (88,1) (90,1) (92,1) (94,1) (96,1) (98,1) (100,1) (102,1) (104,1) (106,1) (108,1) (110,1) (112,1) (114,1) (116,1) (118,1) (120,1) (122,1) (124,1) (126,1) (128,1) (130,1) (132,1) (134,1) (136,1) (138,1) (140,1) (142,1) (144,1) (146,1) (148,1) (150,1) (152,1) (154,1) (156,1) (158,1) (160,1) (162,1) (164,1) (166,1) (168,1) (170,1) (172,1) (174,1) (176,1) (178,1) (180,1) (182,1) (184,1) (186,1) (188,1) (190,1) (192,1) (194,1) (196,1) (198,1) (200,1) (202,1) (204,1) (206,1) (208,1) (210,1) (212,1) (214,1) (216,1) (218,1) (220,1) (222,1) (224,1) (226,1) (228,1) (230,1) (232,1) (234,1) (236,1) (238,1) (240,1) (242,1) (244,1) (246,1) (248,1) (250,1) (252,1) (254,1) (256,1) (258,1) (260,1) (262,1) (264,1) (266,1) (268,1) (270,1) (272,1) (274,1) (276,1) (278,1) (280,1) (282,1) (284,1) (286,1) (288,1) (290,1) (292,1) (294,1) (296,1) (298,1) (300,1)};
\addplot[cgreen,thin,dashed,forget plot] coordinates {(0,0) (2,0) (4,0) (6,0) (8,0) (10,0) (12,0) (14,0) (16,0.066667) (18,0.066667) (20,0.13333) (22,0.13333) (24,0.13333) (26,0.2) (28,0.2) (30,0.31667) (32,0.31667) (34,0.31667) (36,0.38333) (38,0.38333) (40,0.46667) (42,0.46667) (44,0.46667) (46,0.53333) (48,0.53333) (50,0.6) (52,0.6) (54,0.6) (56,0.66667) (58,0.66667) (60,0.73333) (62,0.73333) (64,0.73333) (66,0.8) (68,0.8) (70,0.8) (72,0.8) (74,0.8) (76,0.8) (78,0.8) (80,0.85) (82,0.85) (84,0.85) (86,0.93333) (88,0.93333) (90,0.98333) (92,0.98333) (94,0.98333) (96,0.98333) (98,0.98333) (100,0.98333) (102,0.98333) (104,0.98333) (106,0.98333) (108,0.98333) (110,0.98333) (112,0.98333) (114,0.98333) (116,1) (118,1) (120,1) (122,1) (124,1) (126,1) (128,1) (130,1) (132,1) (134,1) (136,1) (138,1) (140,1) (142,1) (144,1) (146,1) (148,1) (150,1) (152,1) (154,1) (156,1) (158,1) (160,1) (162,1) (164,1) (166,1) (168,1) (170,1) (172,1) (174,1) (176,1) (178,1) (180,1) (182,1) (184,1) (186,1) (188,1) (190,1) (192,1) (194,1) (196,1) (198,1) (200,1) (202,1) (204,1) (206,1) (208,1) (210,1) (212,1) (214,1) (216,1) (218,1) (220,1) (222,1) (224,1) (226,1) (228,1) (230,1) (232,1) (234,1) (236,1) (238,1) (240,1) (242,1) (244,1) (246,1) (248,1) (250,1) (252,1) (254,1) (256,1) (258,1) (260,1) (262,1) (264,1) (266,1) (268,1) (270,1) (272,1) (274,1) (276,1) (278,1) (280,1) (282,1) (284,1) (286,1) (288,1) (290,1) (292,1) (294,1) (296,1) (298,1) (300,1)};
\addplot[cviolet,very thick,forget plot] coordinates {(0,0) (2,0) (4,0) (6,0) (8,0) (10,0) (12,0) (14,0) (16,0.13333) (18,0.13333) (20,0.2) (22,0.2) (24,0.2) (26,0.33333) (28,0.33333) (30,0.4) (32,0.4) (34,0.4) (36,0.46667) (38,0.46667) (40,0.53333) (42,0.53333) (44,0.53333) (46,0.6) (48,0.6) (50,0.66667) (52,0.66667) (54,0.66667) (56,0.76667) (58,0.76667) (60,0.83333) (62,0.83333) (64,0.83333) (66,0.86667) (68,0.86667) (70,0.9) (72,0.9) (74,0.9) (76,0.93333) (78,0.93333) (80,1) (82,1) (84,1) (86,1) (88,1) (90,1) (92,1) (94,1) (96,1) (98,1) (100,1) (102,1) (104,1) (106,1) (108,1) (110,1) (112,1) (114,1) (116,1) (118,1) (120,1) (122,1) (124,1) (126,1) (128,1) (130,1) (132,1) (134,1) (136,1) (138,1) (140,1) (142,1) (144,1) (146,1) (148,1) (150,1) (152,1) (154,1) (156,1) (158,1) (160,1) (162,1) (164,1) (166,1) (168,1) (170,1) (172,1) (174,1) (176,1) (178,1) (180,1) (182,1) (184,1) (186,1) (188,1) (190,1) (192,1) (194,1) (196,1) (198,1) (200,1) (202,1) (204,1) (206,1) (208,1) (210,1) (212,1) (214,1) (216,1) (218,1) (220,1) (222,1) (224,1) (226,1) (228,1) (230,1) (232,1) (234,1) (236,1) (238,1) (240,1) (242,1) (244,1) (246,1) (248,1) (250,1) (252,1) (254,1) (256,1) (258,1) (260,1) (262,1) (264,1) (266,1) (268,1) (270,1) (272,1) (274,1) (276,1) (278,1) (280,1) (282,1) (284,1) (286,1) (288,1) (290,1) (292,1) (294,1) (296,1) (298,1) (300,1)};
\addplot[cviolet,thin,dashed,forget plot] coordinates {(0,0) (2,0) (4,0) (6,0) (8,0) (10,0) (12,0) (14,0) (16,0.066667) (18,0.066667) (20,0.15) (22,0.15) (24,0.15) (26,0.26667) (28,0.26667) (30,0.33333) (32,0.33333) (34,0.33333) (36,0.4) (38,0.4) (40,0.53333) (42,0.53333) (44,0.53333) (46,0.55) (48,0.55) (50,0.61667) (52,0.61667) (54,0.61667) (56,0.68333) (58,0.68333) (60,0.75) (62,0.75) (64,0.75) (66,0.81667) (68,0.81667) (70,0.86667) (72,0.86667) (74,0.86667) (76,0.93333) (78,0.93333) (80,0.93333) (82,0.93333) (84,0.93333) (86,0.95) (88,0.95) (90,1) (92,1) (94,1) (96,1) (98,1) (100,1) (102,1) (104,1) (106,1) (108,1) (110,1) (112,1) (114,1) (116,1) (118,1) (120,1) (122,1) (124,1) (126,1) (128,1) (130,1) (132,1) (134,1) (136,1) (138,1) (140,1) (142,1) (144,1) (146,1) (148,1) (150,1) (152,1) (154,1) (156,1) (158,1) (160,1) (162,1) (164,1) (166,1) (168,1) (170,1) (172,1) (174,1) (176,1) (178,1) (180,1) (182,1) (184,1) (186,1) (188,1) (190,1) (192,1) (194,1) (196,1) (198,1) (200,1) (202,1) (204,1) (206,1) (208,1) (210,1) (212,1) (214,1) (216,1) (218,1) (220,1) (222,1) (224,1) (226,1) (228,1) (230,1) (232,1) (234,1) (236,1) (238,1) (240,1) (242,1) (244,1) (246,1) (248,1) (250,1) (252,1) (254,1) (256,1) (258,1) (260,1) (262,1) (264,1) (266,1) (268,1) (270,1) (272,1) (274,1) (276,1) (278,1) (280,1) (282,1) (284,1) (286,1) (288,1) (290,1) (292,1) (294,1) (296,1) (298,1) (300,1)};
\end{groupplot}
\end{tikzpicture}

\caption{Fraction du front exact déjà simulée, en fonction du nombre de simulations (médiane ; tirets : premier quartile).}\label{fig:conv-front}
\end{figure}

\begin{figure}[htbp]
\centering
% figure générée par tools/make_figures.py : convergence_hv
\begin{tikzpicture}
\begin{groupplot}[group style={group size=2 by 1,horizontal sep=1.4cm},width=0.5\linewidth,height=6cm,
  xlabel={simulations distinctes},ylabel={hypervolume / hypervolume exact},xmin=0,xmax=300,ymin=0,ymax=1.02,
  title style={font=\small},tick label style={font=\scriptsize},
  legend style={draw=none,font=\scriptsize,at={(1.1,-0.25)},anchor=north},legend columns=3,legend cell align=left]
\nextgroupplot[title={S1 : approche de face (7 profils sur le front)}]
\addlegendimage{cmuted,very thick}\addlegendentry{tirage aléatoire}
\addlegendimage{corange,very thick}\addlegendentry{NSGA-II}
\addlegendimage{cblue,very thick}\addlegendentry{boucle QUBO, solveur exact}
\addlegendimage{cgreen,very thick}\addlegendentry{boucle QUBO, recuit simulé}
\addlegendimage{cviolet,very thick}\addlegendentry{boucle QUBO, QAOA $p=2$}
\addplot[cmuted,very thick,forget plot] coordinates {(0,0) (2,0) (4,0.5917) (6,0.5917) (8,0.84018) (10,0.84018) (12,0.87839) (14,0.87839) (16,0.90179) (18,0.90179) (20,0.93005) (22,0.93005) (24,0.94411) (26,0.94411) (28,0.94411) (30,0.94411) (32,0.94871) (34,0.94871) (36,0.94945) (38,0.94945) (40,0.94955) (42,0.94955) (44,0.94955) (46,0.94955) (48,0.94955) (50,0.94955) (52,0.95018) (54,0.95018) (56,0.95018) (58,0.95018) (60,0.95018) (62,0.95018) (64,0.95018) (66,0.95018) (68,0.95023) (70,0.95023) (72,0.95023) (74,0.95023) (76,0.95062) (78,0.95062) (80,0.95068) (82,0.95068) (84,0.95081) (86,0.95081) (88,0.95081) (90,0.95081) (92,0.95081) (94,0.95081) (96,0.95086) (98,0.95086) (100,0.95086) (102,0.95086) (104,0.9514) (106,0.9514) (108,0.9514) (110,0.9514) (112,0.9514) (114,0.9514) (116,0.9514) (118,0.9514) (120,0.9514) (122,0.9514) (124,0.95154) (126,0.95154) (128,0.95154) (130,0.95154) (132,0.96181) (134,0.96181) (136,0.96181) (138,0.96181) (140,0.97214) (142,0.97214) (144,0.97295) (146,0.97295) (148,0.97331) (150,0.97331) (152,0.97331) (154,0.97331) (156,0.97331) (158,0.97331) (160,0.97331) (162,0.97331) (164,0.97448) (166,0.97448) (168,0.97448) (170,0.97448) (172,0.97473) (174,0.97473) (176,0.97475) (178,0.97475) (180,0.97475) (182,0.97475) (184,0.97475) (186,0.97475) (188,0.97475) (190,0.97475) (192,0.97475) (194,0.97475) (196,0.97476) (198,0.97476) (200,0.97496) (202,0.97496) (204,0.97496) (206,0.97496) (208,0.98732) (210,0.98732) (212,0.98735) (214,0.98735) (216,0.98735) (218,0.98735) (220,0.98735) (222,0.98735) (224,0.98735) (226,0.98735) (228,0.98735) (230,0.98735) (232,0.98735) (234,0.98735) (236,0.99957) (238,0.99957) (240,0.99957) (242,0.99957) (244,0.99957) (246,0.99957) (248,0.99963) (250,0.99963) (252,0.99963) (254,0.99963) (256,0.99963) (258,0.99963) (260,0.99976) (262,0.99976) (264,0.99976) (266,0.99976) (268,0.99976) (270,0.99976) (272,0.99976) (274,0.99976) (276,0.99976) (278,0.99976) (280,0.99976) (282,0.99976) (284,0.99976) (286,0.99976) (288,0.99982) (290,0.99982) (292,0.99982) (294,0.99982) (296,0.99982) (298,0.99982) (300,0.99982)};
\addplot[cmuted,thin,dashed,forget plot] coordinates {(0,0) (2,0) (4,0.4106) (6,0.4106) (8,0.74515) (10,0.74515) (12,0.82073) (14,0.82073) (16,0.84066) (18,0.84066) (20,0.86079) (22,0.86079) (24,0.86677) (26,0.86677) (28,0.88702) (30,0.88702) (32,0.9058) (34,0.9058) (36,0.92079) (38,0.92079) (40,0.92079) (42,0.92079) (44,0.92079) (46,0.92079) (48,0.92285) (50,0.92285) (52,0.92285) (54,0.92285) (56,0.92295) (58,0.92295) (60,0.92308) (62,0.92308) (64,0.92308) (66,0.92308) (68,0.92308) (70,0.92308) (72,0.92308) (74,0.92308) (76,0.94049) (78,0.94049) (80,0.94461) (82,0.94461) (84,0.94461) (86,0.94461) (88,0.94463) (90,0.94463) (92,0.94463) (94,0.94463) (96,0.94463) (98,0.94463) (100,0.94463) (102,0.94463) (104,0.94808) (106,0.94808) (108,0.94824) (110,0.94824) (112,0.94824) (114,0.94824) (116,0.94824) (118,0.94824) (120,0.94824) (122,0.94824) (124,0.94962) (126,0.94962) (128,0.94962) (130,0.94962) (132,0.94962) (134,0.94962) (136,0.94962) (138,0.94962) (140,0.94962) (142,0.94962) (144,0.94962) (146,0.94962) (148,0.9503) (150,0.9503) (152,0.9503) (154,0.9503) (156,0.9503) (158,0.9503) (160,0.9503) (162,0.9503) (164,0.95095) (166,0.95095) (168,0.95095) (170,0.95095) (172,0.9513) (174,0.9513) (176,0.9513) (178,0.9513) (180,0.9513) (182,0.9513) (184,0.95151) (186,0.95151) (188,0.95151) (190,0.95151) (192,0.95151) (194,0.95151) (196,0.95151) (198,0.95151) (200,0.96683) (202,0.96683) (204,0.96683) (206,0.96683) (208,0.96683) (210,0.96683) (212,0.96683) (214,0.96683) (216,0.96683) (218,0.96683) (220,0.96683) (222,0.96683) (224,0.97382) (226,0.97382) (228,0.97382) (230,0.97382) (232,0.97382) (234,0.97382) (236,0.97382) (238,0.97382) (240,0.97382) (242,0.97382) (244,0.97382) (246,0.97382) (248,0.97382) (250,0.97382) (252,0.97382) (254,0.97382) (256,0.97388) (258,0.97388) (260,0.97389) (262,0.97389) (264,0.97389) (266,0.97389) (268,0.97389) (270,0.97389) (272,0.97389) (274,0.97389) (276,0.97392) (278,0.97392) (280,0.97392) (282,0.97392) (284,0.97392) (286,0.97392) (288,0.97392) (290,0.97392) (292,0.97392) (294,0.97392) (296,0.97392) (298,0.97392) (300,0.97392)};
\addplot[corange,very thick,forget plot] coordinates {(0,0) (2,0) (4,0) (6,0) (8,0) (10,0) (12,0) (14,0) (16,0.81908) (18,0.81908) (20,0.81908) (22,0.81908) (24,0.81908) (26,0.81908) (28,0.81908) (30,0.81908) (32,0.92577) (34,0.92577) (36,0.92577) (38,0.92577) (40,0.92577) (42,0.92577) (44,0.92577) (46,0.92577) (48,0.9511) (50,0.9511) (52,0.9511) (54,0.9511) (56,0.9511) (58,0.9511) (60,0.9511) (62,0.96131) (64,0.97219) (66,0.97219) (68,0.97219) (70,0.97219) (72,0.97219) (74,0.97219) (76,0.9722) (78,0.97367) (80,0.97367) (82,0.97367) (84,0.97367) (86,0.97514) (88,0.97514) (90,0.97514) (92,0.97517) (94,0.97517) (96,0.99991) (98,0.99999) (100,0.99999) (102,0.99999) (104,1) (106,1) (108,1) (110,1) (112,1) (114,1) (116,1) (118,1) (120,1) (122,1) (124,1) (126,1) (128,1) (130,1) (132,1) (134,1) (136,1) (138,1) (140,1) (142,1) (144,1) (146,1) (148,1) (150,1) (152,1) (154,1) (156,1) (158,1) (160,1) (162,1) (164,1) (166,1) (168,1) (170,1) (172,1) (174,1) (176,1) (178,1) (180,1) (182,1) (184,1) (186,1) (188,1) (190,1) (192,1) (194,1) (196,1) (198,1) (200,1) (202,1) (204,1) (206,1) (208,1) (210,1) (212,1) (214,1) (216,1) (218,1) (220,1) (222,1) (224,1) (226,1) (228,1) (230,1) (232,1) (234,1) (236,1) (238,1) (240,1) (242,1) (244,1) (246,1) (248,1) (250,1) (252,1) (254,1) (256,1) (258,1) (260,1) (262,1) (264,1) (266,1) (268,1) (270,1) (272,1) (274,1) (276,1) (278,1) (280,1) (282,1) (284,1) (286,1) (288,1) (290,1) (292,1) (294,1) (296,1) (298,1) (300,1)};
\addplot[corange,thin,dashed,forget plot] coordinates {(0,0) (2,0) (4,0) (6,0) (8,0) (10,0) (12,0) (14,0) (16,0.78505) (18,0.78505) (20,0.78505) (22,0.78505) (24,0.78505) (26,0.78505) (28,0.78505) (30,0.78505) (32,0.83978) (34,0.83978) (36,0.83978) (38,0.83978) (40,0.83978) (42,0.83978) (44,0.83978) (46,0.83978) (48,0.94577) (50,0.94577) (52,0.94577) (54,0.94577) (56,0.94577) (58,0.94577) (60,0.94577) (62,0.94777) (64,0.9507) (66,0.9507) (68,0.9507) (70,0.9507) (72,0.9507) (74,0.9507) (76,0.96596) (78,0.96677) (80,0.96677) (82,0.96677) (84,0.96677) (86,0.96677) (88,0.96698) (90,0.96707) (92,0.97223) (94,0.97223) (96,0.97223) (98,0.97223) (100,0.97518) (102,0.97518) (104,0.99367) (106,0.99376) (108,0.99376) (110,1) (112,1) (114,1) (116,1) (118,1) (120,1) (122,1) (124,1) (126,1) (128,1) (130,1) (132,1) (134,1) (136,1) (138,1) (140,1) (142,1) (144,1) (146,1) (148,1) (150,1) (152,1) (154,1) (156,1) (158,1) (160,1) (162,1) (164,1) (166,1) (168,1) (170,1) (172,1) (174,1) (176,1) (178,1) (180,1) (182,1) (184,1) (186,1) (188,1) (190,1) (192,1) (194,1) (196,1) (198,1) (200,1) (202,1) (204,1) (206,1) (208,1) (210,1) (212,1) (214,1) (216,1) (218,1) (220,1) (222,1) (224,1) (226,1) (228,1) (230,1) (232,1) (234,1) (236,1) (238,1) (240,1) (242,1) (244,1) (246,1) (248,1) (250,1) (252,1) (254,1) (256,1) (258,1) (260,1) (262,1) (264,1) (266,1) (268,1) (270,1) (272,1) (274,1) (276,1) (278,1) (280,1) (282,1) (284,1) (286,1) (288,1) (290,1) (292,1) (294,1) (296,1) (298,1) (300,1)};
\addplot[cblue,very thick,forget plot] coordinates {(0,0) (2,0) (4,0) (6,0) (8,0) (10,0) (12,0) (14,0) (16,0.92534) (18,0.92534) (20,0.95137) (22,0.95137) (24,0.95137) (26,0.97371) (28,0.97371) (30,0.97518) (32,0.97518) (34,0.97518) (36,1) (38,1) (40,1) (42,1) (44,1) (46,1) (48,1) (50,1) (52,1) (54,1) (56,1) (58,1) (60,1) (62,1) (64,1) (66,1) (68,1) (70,1) (72,1) (74,1) (76,1) (78,1) (80,1) (82,1) (84,1) (86,1) (88,1) (90,1) (92,1) (94,1) (96,1) (98,1) (100,1) (102,1) (104,1) (106,1) (108,1) (110,1) (112,1) (114,1) (116,1) (118,1) (120,1) (122,1) (124,1) (126,1) (128,1) (130,1) (132,1) (134,1) (136,1) (138,1) (140,1) (142,1) (144,1) (146,1) (148,1) (150,1) (152,1) (154,1) (156,1) (158,1) (160,1) (162,1) (164,1) (166,1) (168,1) (170,1) (172,1) (174,1) (176,1) (178,1) (180,1) (182,1) (184,1) (186,1) (188,1) (190,1) (192,1) (194,1) (196,1) (198,1) (200,1) (202,1) (204,1) (206,1) (208,1) (210,1) (212,1) (214,1) (216,1) (218,1) (220,1) (222,1) (224,1) (226,1) (228,1) (230,1) (232,1) (234,1) (236,1) (238,1) (240,1) (242,1) (244,1) (246,1) (248,1) (250,1) (252,1) (254,1) (256,1) (258,1) (260,1) (262,1) (264,1) (266,1) (268,1) (270,1) (272,1) (274,1) (276,1) (278,1) (280,1) (282,1) (284,1) (286,1) (288,1) (290,1) (292,1) (294,1) (296,1) (298,1) (300,1)};
\addplot[cblue,thin,dashed,forget plot] coordinates {(0,0) (2,0) (4,0) (6,0) (8,0) (10,0) (12,0) (14,0) (16,0.92338) (18,0.92338) (20,0.94951) (22,0.94951) (24,0.94951) (26,0.95155) (28,0.95155) (30,0.97223) (32,0.97223) (34,0.97223) (36,0.97518) (38,0.97518) (40,1) (42,1) (44,1) (46,1) (48,1) (50,1) (52,1) (54,1) (56,1) (58,1) (60,1) (62,1) (64,1) (66,1) (68,1) (70,1) (72,1) (74,1) (76,1) (78,1) (80,1) (82,1) (84,1) (86,1) (88,1) (90,1) (92,1) (94,1) (96,1) (98,1) (100,1) (102,1) (104,1) (106,1) (108,1) (110,1) (112,1) (114,1) (116,1) (118,1) (120,1) (122,1) (124,1) (126,1) (128,1) (130,1) (132,1) (134,1) (136,1) (138,1) (140,1) (142,1) (144,1) (146,1) (148,1) (150,1) (152,1) (154,1) (156,1) (158,1) (160,1) (162,1) (164,1) (166,1) (168,1) (170,1) (172,1) (174,1) (176,1) (178,1) (180,1) (182,1) (184,1) (186,1) (188,1) (190,1) (192,1) (194,1) (196,1) (198,1) (200,1) (202,1) (204,1) (206,1) (208,1) (210,1) (212,1) (214,1) (216,1) (218,1) (220,1) (222,1) (224,1) (226,1) (228,1) (230,1) (232,1) (234,1) (236,1) (238,1) (240,1) (242,1) (244,1) (246,1) (248,1) (250,1) (252,1) (254,1) (256,1) (258,1) (260,1) (262,1) (264,1) (266,1) (268,1) (270,1) (272,1) (274,1) (276,1) (278,1) (280,1) (282,1) (284,1) (286,1) (288,1) (290,1) (292,1) (294,1) (296,1) (298,1) (300,1)};
\addplot[cgreen,very thick,forget plot] coordinates {(0,0) (2,0) (4,0) (6,0) (8,0) (10,0) (12,0) (14,0) (16,0.94585) (18,0.94585) (20,0.97219) (22,0.97219) (24,0.97219) (26,0.97492) (28,0.97492) (30,0.99984) (32,0.99984) (34,0.99984) (36,0.99993) (38,0.99993) (40,1) (42,1) (44,1) (46,1) (48,1) (50,1) (52,1) (54,1) (56,1) (58,1) (60,1) (62,1) (64,1) (66,1) (68,1) (70,1) (72,1) (74,1) (76,1) (78,1) (80,1) (82,1) (84,1) (86,1) (88,1) (90,1) (92,1) (94,1) (96,1) (98,1) (100,1) (102,1) (104,1) (106,1) (108,1) (110,1) (112,1) (114,1) (116,1) (118,1) (120,1) (122,1) (124,1) (126,1) (128,1) (130,1) (132,1) (134,1) (136,1) (138,1) (140,1) (142,1) (144,1) (146,1) (148,1) (150,1) (152,1) (154,1) (156,1) (158,1) (160,1) (162,1) (164,1) (166,1) (168,1) (170,1) (172,1) (174,1) (176,1) (178,1) (180,1) (182,1) (184,1) (186,1) (188,1) (190,1) (192,1) (194,1) (196,1) (198,1) (200,1) (202,1) (204,1) (206,1) (208,1) (210,1) (212,1) (214,1) (216,1) (218,1) (220,1) (222,1) (224,1) (226,1) (228,1) (230,1) (232,1) (234,1) (236,1) (238,1) (240,1) (242,1) (244,1) (246,1) (248,1) (250,1) (252,1) (254,1) (256,1) (258,1) (260,1) (262,1) (264,1) (266,1) (268,1) (270,1) (272,1) (274,1) (276,1) (278,1) (280,1) (282,1) (284,1) (286,1) (288,1) (290,1) (292,1) (294,1) (296,1) (298,1) (300,1)};
\addplot[cgreen,thin,dashed,forget plot] coordinates {(0,0) (2,0) (4,0) (6,0) (8,0) (10,0) (12,0) (14,0) (16,0.91849) (18,0.91849) (20,0.95067) (22,0.95067) (24,0.95067) (26,0.97219) (28,0.97219) (30,0.97441) (32,0.97441) (34,0.97441) (36,0.97518) (38,0.97518) (40,0.99986) (42,0.99986) (44,0.99986) (46,1) (48,1) (50,1) (52,1) (54,1) (56,1) (58,1) (60,1) (62,1) (64,1) (66,1) (68,1) (70,1) (72,1) (74,1) (76,1) (78,1) (80,1) (82,1) (84,1) (86,1) (88,1) (90,1) (92,1) (94,1) (96,1) (98,1) (100,1) (102,1) (104,1) (106,1) (108,1) (110,1) (112,1) (114,1) (116,1) (118,1) (120,1) (122,1) (124,1) (126,1) (128,1) (130,1) (132,1) (134,1) (136,1) (138,1) (140,1) (142,1) (144,1) (146,1) (148,1) (150,1) (152,1) (154,1) (156,1) (158,1) (160,1) (162,1) (164,1) (166,1) (168,1) (170,1) (172,1) (174,1) (176,1) (178,1) (180,1) (182,1) (184,1) (186,1) (188,1) (190,1) (192,1) (194,1) (196,1) (198,1) (200,1) (202,1) (204,1) (206,1) (208,1) (210,1) (212,1) (214,1) (216,1) (218,1) (220,1) (222,1) (224,1) (226,1) (228,1) (230,1) (232,1) (234,1) (236,1) (238,1) (240,1) (242,1) (244,1) (246,1) (248,1) (250,1) (252,1) (254,1) (256,1) (258,1) (260,1) (262,1) (264,1) (266,1) (268,1) (270,1) (272,1) (274,1) (276,1) (278,1) (280,1) (282,1) (284,1) (286,1) (288,1) (290,1) (292,1) (294,1) (296,1) (298,1) (300,1)};
\addplot[cviolet,very thick,forget plot] coordinates {(0,0) (2,0) (4,0) (6,0) (8,0) (10,0) (12,0) (14,0) (16,0.94609) (18,0.94609) (20,0.9721) (22,0.9721) (24,0.9721) (26,0.97509) (28,0.97509) (30,0.99999) (32,0.99999) (34,0.99999) (36,1) (38,1) (40,1) (42,1) (44,1) (46,1) (48,1) (50,1) (52,1) (54,1) (56,1) (58,1) (60,1) (62,1) (64,1) (66,1) (68,1) (70,1) (72,1) (74,1) (76,1) (78,1) (80,1) (82,1) (84,1) (86,1) (88,1) (90,1) (92,1) (94,1) (96,1) (98,1) (100,1) (102,1) (104,1) (106,1) (108,1) (110,1) (112,1) (114,1) (116,1) (118,1) (120,1) (122,1) (124,1) (126,1) (128,1) (130,1) (132,1) (134,1) (136,1) (138,1) (140,1) (142,1) (144,1) (146,1) (148,1) (150,1) (152,1) (154,1) (156,1) (158,1) (160,1) (162,1) (164,1) (166,1) (168,1) (170,1) (172,1) (174,1) (176,1) (178,1) (180,1) (182,1) (184,1) (186,1) (188,1) (190,1) (192,1) (194,1) (196,1) (198,1) (200,1) (202,1) (204,1) (206,1) (208,1) (210,1) (212,1) (214,1) (216,1) (218,1) (220,1) (222,1) (224,1) (226,1) (228,1) (230,1) (232,1) (234,1) (236,1) (238,1) (240,1) (242,1) (244,1) (246,1) (248,1) (250,1) (252,1) (254,1) (256,1) (258,1) (260,1) (262,1) (264,1) (266,1) (268,1) (270,1) (272,1) (274,1) (276,1) (278,1) (280,1) (282,1) (284,1) (286,1) (288,1) (290,1) (292,1) (294,1) (296,1) (298,1) (300,1)};
\addplot[cviolet,thin,dashed,forget plot] coordinates {(0,0) (2,0) (4,0) (6,0) (8,0) (10,0) (12,0) (14,0) (16,0.92727) (18,0.92727) (20,0.95067) (22,0.95067) (24,0.95067) (26,0.9729) (28,0.9729) (30,0.98134) (32,0.98134) (34,0.98134) (36,1) (38,1) (40,1) (42,1) (44,1) (46,1) (48,1) (50,1) (52,1) (54,1) (56,1) (58,1) (60,1) (62,1) (64,1) (66,1) (68,1) (70,1) (72,1) (74,1) (76,1) (78,1) (80,1) (82,1) (84,1) (86,1) (88,1) (90,1) (92,1) (94,1) (96,1) (98,1) (100,1) (102,1) (104,1) (106,1) (108,1) (110,1) (112,1) (114,1) (116,1) (118,1) (120,1) (122,1) (124,1) (126,1) (128,1) (130,1) (132,1) (134,1) (136,1) (138,1) (140,1) (142,1) (144,1) (146,1) (148,1) (150,1) (152,1) (154,1) (156,1) (158,1) (160,1) (162,1) (164,1) (166,1) (168,1) (170,1) (172,1) (174,1) (176,1) (178,1) (180,1) (182,1) (184,1) (186,1) (188,1) (190,1) (192,1) (194,1) (196,1) (198,1) (200,1) (202,1) (204,1) (206,1) (208,1) (210,1) (212,1) (214,1) (216,1) (218,1) (220,1) (222,1) (224,1) (226,1) (228,1) (230,1) (232,1) (234,1) (236,1) (238,1) (240,1) (242,1) (244,1) (246,1) (248,1) (250,1) (252,1) (254,1) (256,1) (258,1) (260,1) (262,1) (264,1) (266,1) (268,1) (270,1) (272,1) (274,1) (276,1) (278,1) (280,1) (282,1) (284,1) (286,1) (288,1) (290,1) (292,1) (294,1) (296,1) (298,1) (300,1)};
\nextgroupplot[title={S2 : radar en contrebas (15 profils sur le front)}]
\addplot[cmuted,very thick,forget plot] coordinates {(0,0) (2,0) (4,0.60126) (6,0.60126) (8,0.74469) (10,0.74469) (12,0.78239) (14,0.78239) (16,0.81628) (18,0.81628) (20,0.83319) (22,0.83319) (24,0.8443) (26,0.8443) (28,0.8443) (30,0.8443) (32,0.85021) (34,0.85021) (36,0.85453) (38,0.85453) (40,0.85453) (42,0.85453) (44,0.88644) (46,0.88644) (48,0.88644) (50,0.88644) (52,0.90556) (54,0.90556) (56,0.90556) (58,0.90556) (60,0.90671) (62,0.90671) (64,0.91017) (66,0.91017) (68,0.91115) (70,0.91115) (72,0.91146) (74,0.91146) (76,0.9161) (78,0.9161) (80,0.91924) (82,0.91924) (84,0.91924) (86,0.91924) (88,0.91924) (90,0.91924) (92,0.91924) (94,0.91924) (96,0.91932) (98,0.91932) (100,0.92682) (102,0.92682) (104,0.93674) (106,0.93674) (108,0.9369) (110,0.9369) (112,0.9369) (114,0.9369) (116,0.9369) (118,0.9369) (120,0.9369) (122,0.9369) (124,0.93763) (126,0.93763) (128,0.93797) (130,0.93797) (132,0.94164) (134,0.94164) (136,0.94164) (138,0.94164) (140,0.94164) (142,0.94164) (144,0.94468) (146,0.94468) (148,0.94468) (150,0.94468) (152,0.94468) (154,0.94468) (156,0.94468) (158,0.94468) (160,0.94468) (162,0.94468) (164,0.94863) (166,0.94863) (168,0.94867) (170,0.94867) (172,0.95112) (174,0.95112) (176,0.95876) (178,0.95876) (180,0.95876) (182,0.95876) (184,0.96008) (186,0.96008) (188,0.96012) (190,0.96012) (192,0.96012) (194,0.96012) (196,0.96012) (198,0.96012) (200,0.97039) (202,0.97039) (204,0.97039) (206,0.97039) (208,0.97039) (210,0.97039) (212,0.97121) (214,0.97121) (216,0.97121) (218,0.97121) (220,0.97121) (222,0.97121) (224,0.97121) (226,0.97121) (228,0.97121) (230,0.97121) (232,0.97121) (234,0.97121) (236,0.97121) (238,0.97121) (240,0.97121) (242,0.97121) (244,0.97142) (246,0.97142) (248,0.97162) (250,0.97162) (252,0.97793) (254,0.97793) (256,0.97793) (258,0.97793) (260,0.988) (262,0.988) (264,0.98801) (266,0.98801) (268,0.98801) (270,0.98801) (272,0.98801) (274,0.98801) (276,0.98801) (278,0.98801) (280,0.98801) (282,0.98801) (284,0.98801) (286,0.98801) (288,0.99154) (290,0.99154) (292,0.99215) (294,0.99215) (296,0.99215) (298,0.99215) (300,0.99215)};
\addplot[cmuted,thin,dashed,forget plot] coordinates {(0,0) (2,0) (4,0.40877) (6,0.40877) (8,0.68934) (10,0.68934) (12,0.70437) (14,0.70437) (16,0.73949) (18,0.73949) (20,0.75131) (22,0.75131) (24,0.79561) (26,0.79561) (28,0.81944) (30,0.81944) (32,0.83158) (34,0.83158) (36,0.83158) (38,0.83158) (40,0.83158) (42,0.83158) (44,0.83158) (46,0.83158) (48,0.83627) (50,0.83627) (52,0.84802) (54,0.84802) (56,0.85102) (58,0.85102) (60,0.85102) (62,0.85102) (64,0.85102) (66,0.85102) (68,0.85378) (70,0.85378) (72,0.86471) (74,0.86471) (76,0.87044) (78,0.87044) (80,0.89063) (82,0.89063) (84,0.89472) (86,0.89472) (88,0.89725) (90,0.89725) (92,0.89725) (94,0.89725) (96,0.9034) (98,0.9034) (100,0.90549) (102,0.90549) (104,0.91256) (106,0.91256) (108,0.913) (110,0.913) (112,0.91482) (114,0.91482) (116,0.9156) (118,0.9156) (120,0.91578) (122,0.91578) (124,0.91986) (126,0.91986) (128,0.9204) (130,0.9204) (132,0.9204) (134,0.9204) (136,0.9204) (138,0.9204) (140,0.9221) (142,0.9221) (144,0.92685) (146,0.92685) (148,0.92745) (150,0.92745) (152,0.92762) (154,0.92762) (156,0.92762) (158,0.92762) (160,0.92762) (162,0.92762) (164,0.93329) (166,0.93329) (168,0.93389) (170,0.93389) (172,0.93895) (174,0.93895) (176,0.93895) (178,0.93895) (180,0.93895) (182,0.93895) (184,0.94525) (186,0.94525) (188,0.94525) (190,0.94525) (192,0.94525) (194,0.94525) (196,0.94525) (198,0.94525) (200,0.94525) (202,0.94525) (204,0.94525) (206,0.94525) (208,0.94525) (210,0.94525) (212,0.94525) (214,0.94525) (216,0.94525) (218,0.94525) (220,0.94525) (222,0.94525) (224,0.94612) (226,0.94612) (228,0.94612) (230,0.94612) (232,0.9509) (234,0.9509) (236,0.95434) (238,0.95434) (240,0.95434) (242,0.95434) (244,0.95993) (246,0.95993) (248,0.95997) (250,0.95997) (252,0.95997) (254,0.95997) (256,0.96409) (258,0.96409) (260,0.9649) (262,0.9649) (264,0.9649) (266,0.9649) (268,0.9649) (270,0.9649) (272,0.9649) (274,0.9649) (276,0.9649) (278,0.9649) (280,0.9649) (282,0.9649) (284,0.966) (286,0.966) (288,0.96602) (290,0.96602) (292,0.96631) (294,0.96631) (296,0.96633) (298,0.96633) (300,0.96633)};
\addplot[corange,very thick,forget plot] coordinates {(0,0) (2,0) (4,0) (6,0) (8,0) (10,0) (12,0) (14,0) (16,0.77142) (18,0.77142) (20,0.77142) (22,0.77142) (24,0.77142) (26,0.77142) (28,0.77142) (30,0.77142) (32,0.84697) (34,0.84697) (36,0.84697) (38,0.84697) (40,0.84697) (42,0.84697) (44,0.84697) (46,0.86522) (48,0.932) (50,0.932) (52,0.932) (54,0.932) (56,0.932) (58,0.932) (60,0.93386) (62,0.94045) (64,0.95548) (66,0.95548) (68,0.95548) (70,0.95548) (72,0.95635) (74,0.95893) (76,0.9624) (78,0.97558) (80,0.97558) (82,0.97748) (84,0.97748) (86,0.97773) (88,0.98508) (90,0.99731) (92,0.99808) (94,0.99818) (96,0.99953) (98,0.9999) (100,0.99994) (102,0.99994) (104,0.99996) (106,0.99999) (108,1) (110,1) (112,1) (114,1) (116,1) (118,1) (120,1) (122,1) (124,1) (126,1) (128,1) (130,1) (132,1) (134,1) (136,1) (138,1) (140,1) (142,1) (144,1) (146,1) (148,1) (150,1) (152,1) (154,1) (156,1) (158,1) (160,1) (162,1) (164,1) (166,1) (168,1) (170,1) (172,1) (174,1) (176,1) (178,1) (180,1) (182,1) (184,1) (186,1) (188,1) (190,1) (192,1) (194,1) (196,1) (198,1) (200,1) (202,1) (204,1) (206,1) (208,1) (210,1) (212,1) (214,1) (216,1) (218,1) (220,1) (222,1) (224,1) (226,1) (228,1) (230,1) (232,1) (234,1) (236,1) (238,1) (240,1) (242,1) (244,1) (246,1) (248,1) (250,1) (252,1) (254,1) (256,1) (258,1) (260,1) (262,1) (264,1) (266,1) (268,1) (270,1) (272,1) (274,1) (276,1) (278,1) (280,1) (282,1) (284,1) (286,1) (288,1) (290,1) (292,1) (294,1) (296,1) (298,1) (300,1)};
\addplot[corange,thin,dashed,forget plot] coordinates {(0,0) (2,0) (4,0) (6,0) (8,0) (10,0) (12,0) (14,0) (16,0.73301) (18,0.73301) (20,0.73301) (22,0.73301) (24,0.73301) (26,0.73301) (28,0.73301) (30,0.73301) (32,0.79655) (34,0.79655) (36,0.79655) (38,0.79655) (40,0.79655) (42,0.79655) (44,0.79655) (46,0.79655) (48,0.87838) (50,0.87838) (52,0.87838) (54,0.87838) (56,0.87838) (58,0.87838) (60,0.89493) (62,0.92771) (64,0.94272) (66,0.94272) (68,0.94272) (70,0.94272) (72,0.94272) (74,0.94272) (76,0.94272) (78,0.95441) (80,0.95441) (82,0.95441) (84,0.95441) (86,0.96137) (88,0.9714) (90,0.9714) (92,0.97677) (94,0.97782) (96,0.97782) (98,0.98879) (100,0.98886) (102,0.99897) (104,0.99899) (106,0.99899) (108,0.99899) (110,0.99982) (112,0.99998) (114,1) (116,1) (118,1) (120,1) (122,1) (124,1) (126,1) (128,1) (130,1) (132,1) (134,1) (136,1) (138,1) (140,1) (142,1) (144,1) (146,1) (148,1) (150,1) (152,1) (154,1) (156,1) (158,1) (160,1) (162,1) (164,1) (166,1) (168,1) (170,1) (172,1) (174,1) (176,1) (178,1) (180,1) (182,1) (184,1) (186,1) (188,1) (190,1) (192,1) (194,1) (196,1) (198,1) (200,1) (202,1) (204,1) (206,1) (208,1) (210,1) (212,1) (214,1) (216,1) (218,1) (220,1) (222,1) (224,1) (226,1) (228,1) (230,1) (232,1) (234,1) (236,1) (238,1) (240,1) (242,1) (244,1) (246,1) (248,1) (250,1) (252,1) (254,1) (256,1) (258,1) (260,1) (262,1) (264,1) (266,1) (268,1) (270,1) (272,1) (274,1) (276,1) (278,1) (280,1) (282,1) (284,1) (286,1) (288,1) (290,1) (292,1) (294,1) (296,1) (298,1) (300,1)};
\addplot[cblue,very thick,forget plot] coordinates {(0,0) (2,0) (4,0) (6,0) (8,0) (10,0) (12,0) (14,0) (16,0.88006) (18,0.88006) (20,0.9398) (22,0.9398) (24,0.9398) (26,0.96817) (28,0.96817) (30,0.9726) (32,0.9726) (34,0.9726) (36,0.98705) (38,0.98705) (40,0.99511) (42,0.99511) (44,0.99511) (46,0.9972) (48,0.9972) (50,0.99851) (52,0.99851) (54,0.99851) (56,0.99935) (58,0.99935) (60,0.99996) (62,0.99996) (64,0.99996) (66,0.99997) (68,0.99997) (70,1) (72,1) (74,1) (76,1) (78,1) (80,1) (82,1) (84,1) (86,1) (88,1) (90,1) (92,1) (94,1) (96,1) (98,1) (100,1) (102,1) (104,1) (106,1) (108,1) (110,1) (112,1) (114,1) (116,1) (118,1) (120,1) (122,1) (124,1) (126,1) (128,1) (130,1) (132,1) (134,1) (136,1) (138,1) (140,1) (142,1) (144,1) (146,1) (148,1) (150,1) (152,1) (154,1) (156,1) (158,1) (160,1) (162,1) (164,1) (166,1) (168,1) (170,1) (172,1) (174,1) (176,1) (178,1) (180,1) (182,1) (184,1) (186,1) (188,1) (190,1) (192,1) (194,1) (196,1) (198,1) (200,1) (202,1) (204,1) (206,1) (208,1) (210,1) (212,1) (214,1) (216,1) (218,1) (220,1) (222,1) (224,1) (226,1) (228,1) (230,1) (232,1) (234,1) (236,1) (238,1) (240,1) (242,1) (244,1) (246,1) (248,1) (250,1) (252,1) (254,1) (256,1) (258,1) (260,1) (262,1) (264,1) (266,1) (268,1) (270,1) (272,1) (274,1) (276,1) (278,1) (280,1) (282,1) (284,1) (286,1) (288,1) (290,1) (292,1) (294,1) (296,1) (298,1) (300,1)};
\addplot[cblue,thin,dashed,forget plot] coordinates {(0,0) (2,0) (4,0) (6,0) (8,0) (10,0) (12,0) (14,0) (16,0.85761) (18,0.85761) (20,0.92675) (22,0.92675) (24,0.92675) (26,0.93978) (28,0.93978) (30,0.96543) (32,0.96543) (34,0.96543) (36,0.97566) (38,0.97566) (40,0.9773) (42,0.9773) (44,0.9773) (46,0.9888) (48,0.9888) (50,0.99545) (52,0.99545) (54,0.99545) (56,0.99793) (58,0.99793) (60,0.99867) (62,0.99867) (64,0.99867) (66,0.99876) (68,0.99876) (70,0.99945) (72,0.99945) (74,0.99945) (76,0.99974) (78,0.99974) (80,0.99996) (82,0.99996) (84,0.99996) (86,1) (88,1) (90,1) (92,1) (94,1) (96,1) (98,1) (100,1) (102,1) (104,1) (106,1) (108,1) (110,1) (112,1) (114,1) (116,1) (118,1) (120,1) (122,1) (124,1) (126,1) (128,1) (130,1) (132,1) (134,1) (136,1) (138,1) (140,1) (142,1) (144,1) (146,1) (148,1) (150,1) (152,1) (154,1) (156,1) (158,1) (160,1) (162,1) (164,1) (166,1) (168,1) (170,1) (172,1) (174,1) (176,1) (178,1) (180,1) (182,1) (184,1) (186,1) (188,1) (190,1) (192,1) (194,1) (196,1) (198,1) (200,1) (202,1) (204,1) (206,1) (208,1) (210,1) (212,1) (214,1) (216,1) (218,1) (220,1) (222,1) (224,1) (226,1) (228,1) (230,1) (232,1) (234,1) (236,1) (238,1) (240,1) (242,1) (244,1) (246,1) (248,1) (250,1) (252,1) (254,1) (256,1) (258,1) (260,1) (262,1) (264,1) (266,1) (268,1) (270,1) (272,1) (274,1) (276,1) (278,1) (280,1) (282,1) (284,1) (286,1) (288,1) (290,1) (292,1) (294,1) (296,1) (298,1) (300,1)};
\addplot[cgreen,very thick,forget plot] coordinates {(0,0) (2,0) (4,0) (6,0) (8,0) (10,0) (12,0) (14,0) (16,0.88324) (18,0.88324) (20,0.93363) (22,0.93363) (24,0.93363) (26,0.95234) (28,0.95234) (30,0.96835) (32,0.96835) (34,0.96835) (36,0.9728) (38,0.9728) (40,0.98708) (42,0.98708) (44,0.98708) (46,0.98722) (48,0.98722) (50,0.99338) (52,0.99338) (54,0.99338) (56,0.99744) (58,0.99744) (60,0.99808) (62,0.99808) (64,0.99808) (66,0.99918) (68,0.99918) (70,0.99964) (72,0.99964) (74,0.99964) (76,0.99964) (78,0.99964) (80,0.99994) (82,0.99994) (84,0.99994) (86,1) (88,1) (90,1) (92,1) (94,1) (96,1) (98,1) (100,1) (102,1) (104,1) (106,1) (108,1) (110,1) (112,1) (114,1) (116,1) (118,1) (120,1) (122,1) (124,1) (126,1) (128,1) (130,1) (132,1) (134,1) (136,1) (138,1) (140,1) (142,1) (144,1) (146,1) (148,1) (150,1) (152,1) (154,1) (156,1) (158,1) (160,1) (162,1) (164,1) (166,1) (168,1) (170,1) (172,1) (174,1) (176,1) (178,1) (180,1) (182,1) (184,1) (186,1) (188,1) (190,1) (192,1) (194,1) (196,1) (198,1) (200,1) (202,1) (204,1) (206,1) (208,1) (210,1) (212,1) (214,1) (216,1) (218,1) (220,1) (222,1) (224,1) (226,1) (228,1) (230,1) (232,1) (234,1) (236,1) (238,1) (240,1) (242,1) (244,1) (246,1) (248,1) (250,1) (252,1) (254,1) (256,1) (258,1) (260,1) (262,1) (264,1) (266,1) (268,1) (270,1) (272,1) (274,1) (276,1) (278,1) (280,1) (282,1) (284,1) (286,1) (288,1) (290,1) (292,1) (294,1) (296,1) (298,1) (300,1)};
\addplot[cgreen,thin,dashed,forget plot] coordinates {(0,0) (2,0) (4,0) (6,0) (8,0) (10,0) (12,0) (14,0) (16,0.82868) (18,0.82868) (20,0.90733) (22,0.90733) (24,0.90733) (26,0.93692) (28,0.93692) (30,0.95672) (32,0.95672) (34,0.95672) (36,0.96693) (38,0.96693) (40,0.97336) (42,0.97336) (44,0.97336) (46,0.97639) (48,0.97639) (50,0.98432) (52,0.98432) (54,0.98432) (56,0.99254) (58,0.99254) (60,0.99255) (62,0.99255) (64,0.99255) (66,0.99255) (68,0.99255) (70,0.99702) (72,0.99702) (74,0.99702) (76,0.99703) (78,0.99703) (80,0.9987) (82,0.9987) (84,0.9987) (86,0.99952) (88,0.99952) (90,0.99999) (92,0.99999) (94,0.99999) (96,0.99999) (98,0.99999) (100,0.99999) (102,0.99999) (104,0.99999) (106,0.99999) (108,0.99999) (110,0.99999) (112,0.99999) (114,0.99999) (116,1) (118,1) (120,1) (122,1) (124,1) (126,1) (128,1) (130,1) (132,1) (134,1) (136,1) (138,1) (140,1) (142,1) (144,1) (146,1) (148,1) (150,1) (152,1) (154,1) (156,1) (158,1) (160,1) (162,1) (164,1) (166,1) (168,1) (170,1) (172,1) (174,1) (176,1) (178,1) (180,1) (182,1) (184,1) (186,1) (188,1) (190,1) (192,1) (194,1) (196,1) (198,1) (200,1) (202,1) (204,1) (206,1) (208,1) (210,1) (212,1) (214,1) (216,1) (218,1) (220,1) (222,1) (224,1) (226,1) (228,1) (230,1) (232,1) (234,1) (236,1) (238,1) (240,1) (242,1) (244,1) (246,1) (248,1) (250,1) (252,1) (254,1) (256,1) (258,1) (260,1) (262,1) (264,1) (266,1) (268,1) (270,1) (272,1) (274,1) (276,1) (278,1) (280,1) (282,1) (284,1) (286,1) (288,1) (290,1) (292,1) (294,1) (296,1) (298,1) (300,1)};
\addplot[cviolet,very thick,forget plot] coordinates {(0,0) (2,0) (4,0) (6,0) (8,0) (10,0) (12,0) (14,0) (16,0.89049) (18,0.89049) (20,0.94259) (22,0.94259) (24,0.94259) (26,0.95821) (28,0.95821) (30,0.97257) (32,0.97257) (34,0.97257) (36,0.99033) (38,0.99033) (40,0.99161) (42,0.99161) (44,0.99161) (46,0.9934) (48,0.9934) (50,0.99373) (52,0.99373) (54,0.99373) (56,0.99685) (58,0.99685) (60,0.99721) (62,0.99721) (64,0.99721) (66,0.99874) (68,0.99874) (70,0.99961) (72,0.99961) (74,0.99961) (76,0.99997) (78,0.99997) (80,1) (82,1) (84,1) (86,1) (88,1) (90,1) (92,1) (94,1) (96,1) (98,1) (100,1) (102,1) (104,1) (106,1) (108,1) (110,1) (112,1) (114,1) (116,1) (118,1) (120,1) (122,1) (124,1) (126,1) (128,1) (130,1) (132,1) (134,1) (136,1) (138,1) (140,1) (142,1) (144,1) (146,1) (148,1) (150,1) (152,1) (154,1) (156,1) (158,1) (160,1) (162,1) (164,1) (166,1) (168,1) (170,1) (172,1) (174,1) (176,1) (178,1) (180,1) (182,1) (184,1) (186,1) (188,1) (190,1) (192,1) (194,1) (196,1) (198,1) (200,1) (202,1) (204,1) (206,1) (208,1) (210,1) (212,1) (214,1) (216,1) (218,1) (220,1) (222,1) (224,1) (226,1) (228,1) (230,1) (232,1) (234,1) (236,1) (238,1) (240,1) (242,1) (244,1) (246,1) (248,1) (250,1) (252,1) (254,1) (256,1) (258,1) (260,1) (262,1) (264,1) (266,1) (268,1) (270,1) (272,1) (274,1) (276,1) (278,1) (280,1) (282,1) (284,1) (286,1) (288,1) (290,1) (292,1) (294,1) (296,1) (298,1) (300,1)};
\addplot[cviolet,thin,dashed,forget plot] coordinates {(0,0) (2,0) (4,0) (6,0) (8,0) (10,0) (12,0) (14,0) (16,0.88021) (18,0.88021) (20,0.93052) (22,0.93052) (24,0.93052) (26,0.94529) (28,0.94529) (30,0.95324) (32,0.95324) (34,0.95324) (36,0.97264) (38,0.97264) (40,0.97681) (42,0.97681) (44,0.97681) (46,0.97757) (48,0.97757) (50,0.98665) (52,0.98665) (54,0.98665) (56,0.99248) (58,0.99248) (60,0.9925) (62,0.9925) (64,0.9925) (66,0.99356) (68,0.99356) (70,0.99365) (72,0.99365) (74,0.99365) (76,0.99948) (78,0.99948) (80,0.99948) (82,0.99948) (84,0.99948) (86,0.99997) (88,0.99997) (90,1) (92,1) (94,1) (96,1) (98,1) (100,1) (102,1) (104,1) (106,1) (108,1) (110,1) (112,1) (114,1) (116,1) (118,1) (120,1) (122,1) (124,1) (126,1) (128,1) (130,1) (132,1) (134,1) (136,1) (138,1) (140,1) (142,1) (144,1) (146,1) (148,1) (150,1) (152,1) (154,1) (156,1) (158,1) (160,1) (162,1) (164,1) (166,1) (168,1) (170,1) (172,1) (174,1) (176,1) (178,1) (180,1) (182,1) (184,1) (186,1) (188,1) (190,1) (192,1) (194,1) (196,1) (198,1) (200,1) (202,1) (204,1) (206,1) (208,1) (210,1) (212,1) (214,1) (216,1) (218,1) (220,1) (222,1) (224,1) (226,1) (228,1) (230,1) (232,1) (234,1) (236,1) (238,1) (240,1) (242,1) (244,1) (246,1) (248,1) (250,1) (252,1) (254,1) (256,1) (258,1) (260,1) (262,1) (264,1) (266,1) (268,1) (270,1) (272,1) (274,1) (276,1) (278,1) (280,1) (282,1) (284,1) (286,1) (288,1) (290,1) (292,1) (294,1) (296,1) (298,1) (300,1)};
\end{groupplot}
\end{tikzpicture}

\caption{Hypervolume relatif des profils simulés. La boucle atteint très vite un front presque optimal ; c'est la complétion des derniers points qui coûte.}\label{fig:conv-hv}
\end{figure}

\section{Le tableau de synthèse}
\begin{center}
\renewcommand{\arraystretch}{1.2}
\begin{tabular}{lcccc}
\toprule
 & \multicolumn{2}{c}{S1 (7 profils)} & \multicolumn{2}{c}{S2 (15 profils)} \\
 & front complet & HV 99~\% & front complet & HV 99~\% \\
\midrule
tirage aléatoire & 454 & 222 & 494 & 274 \\
NSGA-II, $N=16$ & 106,5 & 95 & 106 & 89 \\
NSGA-II, $N=8$ (réglé) & 65 & --- & 81 & --- \\
boucle QUBO, solveur exact & \textbf{35} & 35 & \textbf{67,5} & 40 \\
boucle QUBO, recuit simulé & 40 & 30 & 85\textsuperscript{*} & 50 \\
boucle QUBO, QAOA $p=2$ & \textbf{35} & 30 & 80 & \textbf{37,5} \\
\bottomrule
\end{tabular}

\smallskip
{\footnotesize Médianes en nombre de simulations distinctes. \textsuperscript{*}Une graine sur 20 n'a pas complété le front dans le budget de 150 simulations.}
\end{center}

\section{Lecture}
\paragraph{1. Toutes les méthodes battent le hasard, de loin.} Pour avoir tout le front, le tirage aléatoire doit simuler presque toute la grille (454 et 494 profils sur 512) : les derniers profils du front sont trouvés au hasard parmi les tout derniers tirages.

\paragraph{2. La boucle QUBO économise des simulations, surtout sur le front court.} Dans S1, elle retrouve le front en 35 simulations, trois fois moins que NSGA-II par défaut et presque deux fois moins que NSGA-II bien réglé. Dans S2, l'écart se réduit : 67,5 contre 106 et 81. Face à un NSGA-II réglé, le gain n'est plus que de 17~\% avec le solveur exact, et devient nul avec QAOA.

\paragraph{3. Trouver un bon front est facile ; le compléter coûte.} Pour atteindre 99~\% de l'hypervolume, la boucle n'a besoin que de 30 à 50 simulations dans les deux scénarios, deux fois moins que NSGA-II. La complétion du front de S2 est plus coûteuse parce que 10 de ses 15 profils ne sont pas supportés : la somme pondérée ne les vise jamais directement, et la boucle ne les trouve qu'indirectement (exclusion des profils déjà simulés, erreurs du modèle quadratique).

\paragraph{4. Le solveur quantique ne change presque rien.} Avec 13 qubits, le QUBO est résolu exactement en une fraction de seconde. QAOA ($p=2$) et le recuit renvoient parfois un candidat un peu moins bon que l'optimum du QUBO, ce qui ajoute quelques itérations : $+0$ à $+13$ simulations pour QAOA, $+5$ à $+18$ pour le recuit. \textbf{Tout le gain de la boucle vient du modèle de substitution}, qui exploite la régularité des deux objectifs ; le solveur n'est qu'un composant interchangeable.

\section{Ce que coûterait une vraie machine}
Une itération de la boucle avec QAOA, c'est 5 QUBO (un par poids), chacun optimisé par COBYLA avec 3 départs, soit environ 600 évaluations de l'énergie moyenne. Sur un ordinateur quantique, chaque évaluation est une estimation statistique sur des centaines à des milliers de tirs. Avec environ 14 itérations pour S2 :
\[
14\times5\times600\approx4\cdot10^4\ \text{évaluations}\quad\Rightarrow\quad\text{de l'ordre de }10^7\text{ exécutions de circuit à }1\,000\text{ tirs.}
\]
Chaque circuit compte environ 69 couplages, donc 138 CNOT par couche et 276 pour $p=2$, avant tout routage sur la topologie réelle. C'est faisable, mais coûteux et bruité, pour un problème qu'un ordinateur portable résout exactement en une milliseconde. Le seul argument quantique tenable porte donc sur le passage à l'échelle, et c'est ce que doit mesurer l'itération suivante.

\begin{resultat}
\begin{itemize}
\item La boucle QUBO retrouve le front exact avec 1,2 à 3 fois moins de simulations que NSGA-II, et atteint 99~\% de l'hypervolume deux fois plus vite.
\item Ce gain vient du modèle de substitution ; le solveur (exact, recuit, QAOA) ne change presque rien à 13 qubits.
\item La somme pondérée laisse de côté les points non supportés, qui font l'essentiel du coût de complétion quand le front est long.
\end{itemize}
\end{resultat}

\begin{projet}
C'est le résultat de l'itération V1. Il répond à la question de la note d'intention, mais il montre aussi ses limites : le problème est presque dégénéré et trop petit pour que le quantique joue un rôle. D'où la reformulation de l'itération V2 (voir la fiche projet).
\end{projet}

\exercices
\begin{exercice}
À partir du tableau, calculer le facteur de gain de la boucle (solveur exact) par rapport à NSGA-II réglé ($N=8$), pour chaque scénario. Pourquoi le gain diminue-t-il quand le front s'allonge ?
\end{exercice}
\begin{exercice}
On veut estimer l'énergie moyenne d'un circuit QAOA à $\pm0{,}01$ près, l'énergie normalisée ayant un écart-type de l'ordre de 1. Combien de tirs faut-il par évaluation ? En déduire le nombre total de tirs pour une exécution de la boucle sur S2.
\end{exercice}

\begin{pistes}
\begin{itemize}
\item Mesurer la boucle avec un critère d'arrêt réaliste (stagnation de l'hypervolume) plutôt qu'avec le front exact.
\item Réutiliser les angles optimaux de QAOA d'une itération à l'autre : le nombre d'évaluations diminuerait fortement.
\item Ajouter l'optimisation bayésienne multi-objectif comme concurrent classique de même nature que la boucle.
\end{itemize}
\end{pistes}
